% ==================================================================
%  ProposedMethod.tex  --  Section 3 (title and introduction)
%  Structure of Section 3:
%    3.1 Client-side Preprocessing   (ClientPreprocessing.tex, RoundCount.tex)
%    3.2 Server-side Computation     (NeighborCheck.tex: 3.2 header + 3.2.1,
%                                     CoreCheck.tex 3.2.2,
%                                     LabelPropagation.tex 3.2.3,
%                                     Correctness.tex 3.2.4)
%    3.3 Client-side Post-processing (PostProcessing.tex)
%  Cross-references used here:
%    \label{sec:ckks} \label{sec:comparison} \label{sec:model}
%    \label{sec:notation}                                       (Section 2)
%    \label{sec:client} \label{sec:server} \label{sec:core} \label{sec:lp}
%    \label{sec:correct} \label{sec:post}                       (Section 3)
% ==================================================================
\section{Proposed Method}
\label{sec:method}

Every decision the server makes in this section is a comparison: whether
two points are neighbors, whether a point has enough neighbors, and which
of two labels is larger. Comparison is not among the operations of
Section~\ref{sec:ckks}, so the server evaluates it with the composite
minimax approximation $\mathsf{Sgn}$ of Section~\ref{sec:comparison}. The
approximation is defined only on $[-1,1]$, so each input is divided by a
public constant that bounds its magnitude. The neighbor counts and the
labels are bounded by construction (Sections~\ref{sec:core}
and~\ref{sec:lp}). For the distances, this bound is the reason the client
normalizes the data before encrypting it (Section~\ref{sec:client}).

This section gives the algorithms behind the three server stages of
Fig.~\ref{fig:overview}, Neighbor Check, Core Check and Label
Propagation, together with the client-side work that precedes and
follows them. Section~\ref{sec:client} describes how the client prepares
its data. Section~\ref{sec:server} writes each server stage as an
operation on ciphertexts, in the notation of Section~\ref{sec:notation}.
Neighbor Check decides which pairs of points are neighbors and outputs
the result as mask vectors, which Core Check uses next. Core Check
decides which points are core points and outputs a mask vector, which
Label Propagation uses next to turn the masks into cluster labels.
Section~\ref{sec:correct} examines how
accurate each comparison is, and Section~\ref{sec:post} describes the
post-processing on the client.
